% language=us

% \setupbackend[format=pdf/ua-2]
% \setuptagging[state=start]
% \enabledirectives[math.extensibles=both] % todo, when tagging !

% \nopdfcompression

% \enabletrackers[structures.tags.blobs]
% \enabletrackers[structures.tags]
% \enabletrackers[structures.tags.internals]
% \enabletrackers[structures.tags.showtree]

\usemodule[present-boring,abbreviations-logos]

% \definelayer
%   [extratext]
%   [preset=rightbottom,
%    height=\textheight,
%    width=\textwidth]

% \setupbackgrounds
%   [text]
%   [background=extratext]

% \definehighlight[notabene][style=bf]
% \definehighlight[notabene][style=\glyphweight100\underbar]

% \setuptyping
%   [align=hangright,
%    numbering=]

\startdocument
  [title={Fonts},
   banner={another year of playing},
   location={bachotex\enspace {\bf 2024}\enspace meeting}]

\starttitle[title=How does \TEX\ see them]

\startitemize
\startitem
    The typesetting engine sees nodes: glyphs, rules, boxes, glues, kerns,
    penalties and some more. Only the first three have width, height and depth.
\stopitem
\startitem
    There are also discretionaries that can have various sublists with nodes.
    When the engine looks are the list is has to do into these sub lists: pre,
    post and replace.
\stopitem
\startitem
    The engine itself (we're spreaking \LUAMETATEX\ here) does little with fonts:
    it only needs these dimensions.
\stopitem
\startitem
    Glyphs (coming from characters) can have more properties (italic correction,
    anchors, kerning, math properties). But we might need to fix them, just as
    we might want to add additional features.
\stopitem
\startitem
    We delegate combining glyphs into other glyphs to \LUA, but still we need to
    pass these dimensions to the engine.
\stopitem
\startitem
    So called variable fonts are not that dynamic as we use instances so in
    practice then are not different from regular fonts but we need to assembled
    them.
\stopitem
\startitem
    We can initialize and apply font features beforehand, or we can delay it till
    needed (aka dynamic fonts in context) which ads a bit of overhead but we can
    gain back.
\stopitem
\stopitemize

\stoptitle

\starttitle[title=What does the backend]

\startitemize
\startitem
    The (\PDF) backend turns references to glyphs in fonts into (in our case)
    \PDF\ page stream commands that refer to indices.
\stopitem
\startitem
    It only embeds the subset of glyphs from a font that are really used.
\stopitem
\startitem
    It has to deal with so called virtual glyphs that can refer to the same font
    but also to other fonts.
\stopitem
\startitem
    The backend has to apply properties like slant, scales, weights, expansion,
    rotation etc and also track drift.
\stopitem
\startitem
    It has to support \TYPEONE, \OPENTYPE, \TRUETYPE, \TYPETHREE, and inline
    graphic fonts (although these can be \TYPETHREE).
\stopitem
\startitem
    Color fonts are a special case and can be done with runtime generated
    \TYPETHREE\ fonts.
\stopitem
\startitem
    When needed font resources have to be fixed (like turning \TYPEONE\ \CFF1\
    into \OPENTYPE\ \CFF2).
\stopitem
\stopitemize

\stoptitle

\starttitle[title=Compact font mode]

\startitemize
\startitem
    The engine can scale glyphs on demands (\typ {\glyphscale}, \typ
    {\glyphxscale}, \typ {\glyphyscale}), make them oblique (\typ {\glyphslant})
    and make them bolder (\typ {\glyphweight}).
\stopitem
\startitem
    This permits more lightweight system of definitions (\type {\fontspecdef},
    kind of like glyph spec) which needs less loading but on the other hand more
    processing.
\stopitem
\startitem
    The same is true for math fonts where, as script and script script is built in,
    we can stick to one font and select on the fly. Here we have additional scale
    factors.
\stopitem
\startitem
    It makes \CONTEXT\ runs a little faster, depending on the job, and use less
    memory, especially on huge fonts.
\stopitem
\stopitemize

\startbuffer
test {\glyphxscale 2500 test} \quad {\glyphyscale 2500 test} \quad
{\glyphscale 2500 test} \quad {\glyphscale 2000 \glyphweight 100
\glyphslant 200 test} \quad $x=2$ \quad {\glyphscale 1500 $x=2$}
\stopbuffer

\blank[2*big] \typebuffer

\blank[2*big] \start \getbuffer \stop

\stoptitle

\starttitle[title=Math fonts]

\startitemize
\startitem
    After many experiments we made some decisions with respect to math fonts
    because we assume that what we have now will not change; if so we will
    use the version to discriminate tweaks.
\stopitem
\startitem
    We fix and enhance the font at runtime using so called tweaks, an \MKIV\
    mechanism that has been extended in \LMTX.
\stopitem
\startitem
    We drop italic correction and add it to the width. At the same time we add
    right bottom corner kerns (not to be confused with staircase kerns).
\stopitem
\startitem
    We overload top anchors and also add bottom ones when applicable. They can be
    wrong anyway.
\stopitem
\startitem
    We use companion fonts that have consistent dimensions etc. of fence characters
    like parenthesis and vertical rules. These are shipped with \CONTEXT\ installations.
\stopitem
\startitem
    We fix and extensibles, radicals, integrals etc. We also can fix accents.
\stopitem
\startitem
    We introduce additional alphabets and fine tune problematic characters. What we
    really do can differ per font.
\stopitem
\startitem
    The enhanced math engine handles spacing based on classes so we don't need
    glyphs to have spacing artifacts.
\stopitem
\stopitemize

\starttitle[title=Compact math]

\startitemize
\startitem
    We can use one font instance and select script and script-script alternates
    when needed and then apply a related scale.
\stopitem
\startitem
    As with glyphs we can scale horizontal, vertical, and both. So, we have the
    main scale, three extra scales, and the math scales.
\stopitem
\startitem
    Font parameters are scale on the fly too. It all means more overhead but the
    number of math characters used in a formula is not large.
\stopitem
\startitem
    Boldening can be handled in a dedicated instance, after all tweaks might
    adapt on it.
\stopitem
\startitem
    Glyphs can have offsets, alternate dimensions etc either via virtual copies
    or via glyph parameters for the backend.
\stopitem
\startitem
    In \CONTEXT\ we support bidirectional math directly by borrowing and|/|or
    mirroring shapes.
\stopitem
\stopitemize

\stoptitle

\starttitle[title=How about bitmaps]

\startitemize
\startitem
    In \TEX\ backends bitmap fonts were supported by using \TYPETHREE\ fonts.
\stopitem
\startitem
    We also use \TYPETHREE\ for e.g. color fonts, composed math glyphs etc. That way we
    stay in the \UNICODE\ domain.
\stopitem
\startitem
    The way bitmaps are handles in \TEX\ distributions is somewhat complex and it has
    some hard coded assumptions (related to mode files).
\stopitem
\startitem
    We can revive bitmap fonts and even auto convert them in proper outlines.
\stopitem
\stopitemize

\stoptitle

\starttitle[title=Messing with \PDF\ files ]

\startitemize
\startitem
    The last years several somewhat special \PDF\ features were added to
    \CONTEXT\ like signing and encryption.
\stopitem
\startitem
    Winter 2024 manipulating and fixing \PDF\ files was added to this: we have a
    framework for doing this.
\stopitem
\startitem
    This also concerned adding missing fonts, fixing bad inclusion, adapting the
    page stream, removing rubish and recoloring.
\stopitem
\startitem
    Around that time inclusion of \PDF\ pages was extended to also merge fonts
    when possible. (We already had support for hyperlinks, layers, etc. so it
    fits in the repertoire.)
\stopitem
\stopitemize

\starttitle[title=Merging fonts]

\startitemize
\startitem
    When embedding pages from other files we can try to share fonts with the parent document.
\stopitem
\startitem
    This is possible to some extend but it depends on what is known about the other font.
\stopitem
\startitem
    We can do it more reliable in a \CONTEXT\ only workflow because we can add information
    to the child document. (For instance we also need to handle variable fonts.)
\stopitem
\startitem
    This permits efficient sub run examples in for instance manuals.
\stopitem
\startitem
    Merging \TYPEONE\ (or related bitmaps) depends on the availability of fonts
    on the system and encoding vectors in the \PDF\ file.
\stopitem
\startitem
    Merging \OPENTYPE\ and \TRUETYPE\ depends on the indexing method but for
    simple documents \UNICODE\ information can be used.
\stopitem
\startitem
    These methods will be refined as we go.
\stopitem
\stopitemize

\stoptitle

\stoptitle

\starttitle[title=Goodie files]

What are these math companion fonts.

Maybe we show a lfg file (aka font goodie or math tweaks) .

\stoptitle

% \starttitle[title=Fixing \TYPEONE\ fonts]

% \startitemize
% \startitem
% \stopitem
% \startitem
% \stopitem
% \startitem
% \stopitem
% \stopitemize

% \stoptitle

\stopdocument
